\documentclass[10pt]{article}

\usepackage[utf8]{inputenc}

\usepackage[T1]{fontenc}

\usepackage{amsmath}

\usepackage{amsfonts}

\usepackage{amssymb}

\usepackage[version=4]{mhchem}

\usepackage{stmaryrd}

\usepackage{graphicx}

\usepackage[export]{adjustbox}

\graphicspath{ {./images/} }



\begin{document}

показывает, что с увеличением $x$ она становится в среднем все более редкоӥ. Существуют сколь угодно длинные отрезки ряда натуральных чисел, среди к-рых нет ни одного п. ч. Напр., $n-1$ натуральных чисел вида $n !+2, \ldots, n !+n$ для любого $n \geqslant 2$ являются составными числами. В то же время в (1) встречаются п. ч. такие, как 8004119 и 8004121 , разность между к-рыми равна 2 (п. ч.- близнецы). Проблема поведения $\pi(x)$ при $x \rightarrow \infty$ является одной из наиболее трудных и интересных проблем теории чисел.



Первый результат о величине $\pi(x)$ - т е о р е м а Е в кли д а: $\pi(x) \rightarrow \infty$ при $x \rightarrow \infty$. Л. Эйлер (L. Еuler, 1737, 1749, см. [1]) ввел функцию



\[

\zeta(s)=\sum_{n=1}^{\infty} n^{-s}

\]



и показал, что при $s>1$



\[

\zeta(s)=\coprod_{p}\left(1-\frac{1}{p^{s}}\right)^{-1},

\]



где ряд распространяется на все натуральные числа, а произведение на все п. ч. Тождество (3) и его обобщения играют фундаментальную роль в теории Р. п.ч. Исходя из него, Л. Эйлер доказал, что ряд $\sum^{1 / p}$ и произведение $\prod\left(1-\frac{1}{p}\right)^{-1}$ по простым $p$ расходятся. Этоновое доказательство бесконечности числа п. ч. Более того, JI. Эїлер установил, что п. ч. «много", ибо



\[

\pi(x)>\log \frac{x}{e}

\]



и, в то же время, потти все натуральные числа являются составнымі, т. к.



\[

\pi(x) x^{-1} \longrightarrow 0 \text { при } x \longrightarrow \infty .

\]



Далее значительного успеха достиг П. Л. Чебышев (1851-52, см. [2]). Он доказал, что:



\begin{enumerate}

  \item для любых $m>0, M>0$ есть последовательности $x \rightarrow \infty, y \rightarrow \infty$, для к-рых

\end{enumerate}



\[

\begin{gathered}

\pi(x)-\operatorname{li} x<M x \log ^{-m} x, \\

\pi(y)-\operatorname{li} y>-M y \log ^{-m} y ;

\end{gathered}

\]



\begin{enumerate}

  \setcounter{enumi}{1}

  \item если существует предел частного $\pi(x) \log x / x$ при $x \rightarrow \infty$, то он равен 1 .

\end{enumerate}



Тем самым впервые был решен вопрос о существовании простої функции



\[

\begin{gathered}

\text { li } x \equiv \int_{2}^{x} \frac{d u}{\log u}= \\

=\frac{x}{\log x}+\frac{x}{\log ^{2} x}+\ldots+(r-1) ! \frac{x}{\log ^{r} x}+O\left(\frac{x}{\log ^{r+1} x}\right),

\end{gathered}

\]



к-рая служит наилучшим приближением для $\pi(x)$.



Затем П. Л. Чебышев установил истинный порядок роста $\pi(x)$, т. е. существование постоянных $a>0, A>0$ таких, что



\[

a \frac{x}{\log x}<\pi(x)<A \frac{x}{\log x},

\]



причем, $a=0,92 \ldots, A=1,05 \ldots$ для $x \geqslant x_{0}$. Он же доказал, что прі любом $n \geqslant 2$ в интервале $(n, 2 n)$ содержится по краӥнеӥ мере одно п. ч. (І о с т у л а т Б е р т р а н а). В основе вывода неравенств (4) лежит тождество Ч е бы ш е в



\[

F(x) \equiv \log [x] !=\sum_{n \leqslant x} \psi\left(\frac{x}{n}\right),

\]



в к-ром введенная П. Л. Чебышевым функция $\psi$ определяется суммой по степеням $p^{m}, m=1,2, \ldots$, п. ч. $p$ :



\[

\psi(x)=\sum_{p^{m} \leqslant x} \log p=\sum_{n \leqslant x} \Lambda(n) .

\]



Именно, комбпнация $\log [u]$ ! для $x \geqslant 30$ в форме



\[

F^{*}(x)=F(x)+F\left(\frac{x}{30}\right)-F\left(\frac{x}{2}\right)-F\left(\frac{x}{3}\right)-F\left(\frac{x}{6}\right),

\]



вследствие (5), дает тождество



\[

F^{*}(x)=\psi(x)-\psi\left(\frac{x}{i}\right)+\psi\left(\frac{x}{7}\right)-\psi\left(\frac{x}{10}\right)+\ldots,

\]



из к-рого следует, что



\[

\psi(x)-\psi\left(\frac{x}{6}\right)<F^{*}(x)<\psi(x) .

\]



Отсюда, вследствие асимптотич. формулы Стирлинга для $n$ !, вытекает аналог неравенств (4) для $\psi(x)$, из к-рых частичным суммированием получаются неравенства (4). Функция Чебышева $\psi(x)$ оказалась более удобной, чем $\pi(x)$, при изучении Р. п. ч., поскольку наилучшим приближением ее является сам аргумент $x$. Поәтому обычно сначала рассматривают $\psi(x)$, а затем частичным суммированием получают соответствующий результат для $\pi(x)$.



Принцип Римана. В 1859-60 Б. Риман (B. Riemann, см. [3]) рассмотрел введенную Л. Эйлером пля $s>1$ функцию $\zeta(s)$ как функцию комплексного переменного $s=\sigma+i t$, где $\sigma, t$ - действительные переменные, определяемую рядом (2) при $\boldsymbol{\sigma}_{>1}$ (см. Дзета-функция), и обнаружил исключительную важность этой функции для теории Р. п. ч. В частности, он указал выражение разности $\pi(x)$-li $x$ через $x$ и нули функции $\zeta(s)$, лежащие в полосе $0 \leqslant \sigma \leqslant 1$, к-рые наз. н е т р и в и а л ьны м и в у л я м и функции $\zeta(s)$.



Вместо формулы Римана обычно используется более простой конечный ее аналог для $\psi(x)$, доказанный (наряду с формулой Римана) Х. Мангольдтом (Н. Mangoldt, 1895). Именно, для $x>1$



\[

\psi(x)=x-\sum_{|\nu| \leqslant T} \frac{x^{\rho}}{\rho}+O\left(\frac{x}{T} \log ^{2} x T+\log 2 x\right),(6)

\]



где $\rho=\beta+i \gamma$ пробегает нетривиальные нулш $\zeta(s), T-$ любое $\geqslant 2$.



Поскольку



\[

\sum_{|\gamma| \leqslant T} \frac{1}{\gamma} \ll \log ^{2} T

\]



формула (6) показывает, что величина разности $\psi(x)-x$ в главном определяется величиної $\beta$ (действительной частью самых правых нулей $\rho)$. В частности, если $\zeta(s) \neq$ $\neq 0$ правее вертикали $\sigma=\theta, 1 / 2 \leqslant \theta<1$, для функций $\psi(x), \pi(x)$ справедливы следующие асимптотич. выражения:



\[

\begin{gathered}

\psi(x)=x+O\left(x^{\theta} \log ^{2} x\right), \\

\pi(x)=\operatorname{li} x+O\left(x^{\theta} \log x\right) .

\end{gathered}

\]



Наоборот, из этих соотношений следует, что $\zeta(s) \neq 0$ для $\sigma \geqslant \theta$. Если сшраведлива гипотеза Римана, т. е. все нетривиальные нули $\zeta(s)$ лежат на одной прямой $\sigma=1 / 2$, тогда асимптотич. закон распределения п. ч. должен иметь вид



\[

\begin{gathered}

\psi(x)=x+O\left(\sqrt{x} \log ^{2} x\right), \\

\pi(x)=\operatorname{li} x+O(\sqrt{x} \log x) .

\end{gathered}

\]



Причем эти соотношения существенно усилить нельзя, т. е. существуют последовательности $x \rightarrow \infty, y \rightarrow \infty$ такие, что



\[

\begin{aligned}

& \pi(x)-\operatorname{li} x<-\sqrt{x} \frac{\log \log \log x}{\log \log x},

\end{aligned}

\]



\begin{center}

\includegraphics[max width=\textwidth]{2023_07_08_8359c2a82d81740a618bg-1}

\end{center}



Таким образом, согласно принцииу Римана проблема асимптотич. выражения функциї Р. п. ч. $\psi(x), \pi(x)$ сводится к проблеме границы действительной части нетривиальных нулей функции $\zeta(s)$. До сих пор (1983), однако, не удалось наӥти какое-либо постоянное $\theta$, $1 / 2 \leqslant \theta<1$, с условием $\beta \geqslant 0$. Искомая граница для $\beta$





\end{document}